\documentclass[10pt,a4paper]{amsart}
\usepackage{amsmath,amssymb,amsthm,graphicx}
% Retained preamble for the Stacks Project chapter divisors.tex (Divisors).
% The chapter file pinned-original-source.tex carries only \input{preamble} on its first line, so no
% preamble text lives inside the chapter itself.  The Stacks Project uses one single shared preamble for
% every chapter; the snapshot stages it verbatim in sources/stacks-chow/preamble.tex.  The block below is
% copied mechanically, line for line, from that file's lines 14-18 (the Xy-pic setup) and lines 154-217
% (the native theorem declarations together with the mathematical macros).  Nothing was added, removed or
% redefined.  The wrapper class and the amsmath/amssymb/amsthm/graphicx base packages are supplied by the
% assembler, so the class/\externaldocument part of the shared preamble is not repeated here.
% For commutative diagrams we use Xy-pic
\usepackage[all]{xy}

% We use 2cell for 2-commutative diagrams.
\xyoption{2cell}
% Theorem environments.
%
\theoremstyle{plain}
\newtheorem{theorem}[subsection]{Theorem}
\newtheorem{proposition}[subsection]{Proposition}
\newtheorem{lemma}[subsection]{Lemma}

\theoremstyle{definition}
\newtheorem{definition}[subsection]{Definition}
\newtheorem{example}[subsection]{Example}
\newtheorem{exercise}[subsection]{Exercise}
\newtheorem{situation}[subsection]{Situation}

\theoremstyle{remark}
\newtheorem{remark}[subsection]{Remark}
\newtheorem{remarks}[subsection]{Remarks}

\numberwithin{equation}{subsection}

% Macros
%
\def\lim{\mathop{\mathrm{lim}}\nolimits}
\def\colim{\mathop{\mathrm{colim}}\nolimits}
\def\Spec{\mathop{\mathrm{Spec}}}
\def\Hom{\mathop{\mathrm{Hom}}\nolimits}
\def\Ext{\mathop{\mathrm{Ext}}\nolimits}
\def\SheafHom{\mathop{\mathcal{H}\!\mathit{om}}\nolimits}
\def\SheafExt{\mathop{\mathcal{E}\!\mathit{xt}}\nolimits}
\def\Sch{\mathit{Sch}}
\def\Mor{\mathop{\mathrm{Mor}}\nolimits}
\def\Ob{\mathop{\mathrm{Ob}}\nolimits}
\def\Sh{\mathop{\mathit{Sh}}\nolimits}
\def\NL{\mathop{N\!L}\nolimits}
\def\CH{\mathop{\mathrm{CH}}\nolimits}
\def\proetale{{pro\text{-}\acute{e}tale}}
\def\etale{{\acute{e}tale}}
\def\QCoh{\mathit{QCoh}}
\def\Ker{\mathop{\mathrm{Ker}}}
\def\Im{\mathop{\mathrm{Im}}}
\def\Coker{\mathop{\mathrm{Coker}}}
\def\Coim{\mathop{\mathrm{Coim}}}

% Boxtimes
%
\DeclareMathSymbol{\boxtimes}{\mathbin}{AMSa}{"02}

%
% Macros for moduli stacks/spaces
%
\def\QCohstack{\mathcal{QC}\!\mathit{oh}}
\def\Cohstack{\mathcal{C}\!\mathit{oh}}
\def\Spacesstack{\mathcal{S}\!\mathit{paces}}
\def\Quotfunctor{\mathrm{Quot}}
\def\Hilbfunctor{\mathrm{Hilb}}
\def\Curvesstack{\mathcal{C}\!\mathit{urves}}
\def\Polarizedstack{\mathcal{P}\!\mathit{olarized}}
\def\Complexesstack{\mathcal{C}\!\mathit{omplexes}}
% \Pic is the operator that assigns to X its picard group, usage \Pic(X)
% \Picardstack_{X/B} denotes the Picard stack of X over B
% \Picardfunctor_{X/B} denotes the Picard functor of X over B
\def\Pic{\mathop{\mathrm{Pic}}\nolimits}
\def\Picardstack{\mathcal{P}\!\mathit{ic}}
\def\Picardfunctor{\mathrm{Pic}}
\def\Deformationcategory{\mathcal{D}\!\mathit{ef}}

% --- appended 2026-10-12 (Chief): the project own definitions of the reference/slogan/history
% display environments, copied mechanically from the shared Stacks preamble lines 8-12
% (\usepackage{verbatim} plus the three \newenvironment lines, which the chapter preamble keeps
% verbatim from the project preamble). Nothing added, removed or redefined: these environments
% carry no mathematical content and the project itself typesets them as comments. Without these
% three lines a counted statement whose own text uses them fails with Environment slogan undefined.
\usepackage{verbatim}
\newenvironment{reference}{\comment}{\endcomment}
\newenvironment{slogan}{\comment}{\endcomment}
\newenvironment{history}{\comment}{\endcomment}
\usepackage{cleveref}
\crefname{originallemma}{original lemma}{original lemmas}
\Crefname{originallemma}{Original lemma}{Original lemmas}
\makeatletter
\newlabel{algebra-lemma-formally-etale-etale}{{00UR}{source}{}{sourcepointer-1}{}}
\newlabel{algebra-lemma-formally-etale-etale@cref}{{[originallemma][0][]00UR}{[1][0][]source}{}{}{}}
\makeatother
\crefname{originallemma}{original lemma}{original lemmas}
\Crefname{originallemma}{Original lemma}{Original lemmas}
\makeatletter
\newlabel{algebra-lemma-formally-etale-lift-infinitesimal}{{0H1D}{source}{}{sourcepointer-2}{}}
\newlabel{algebra-lemma-formally-etale-lift-infinitesimal@cref}{{[originallemma][0][]0H1D}{[1][0][]source}{}{}{}}
\makeatother
\crefname{originallemma}{original lemma}{original lemmas}
\Crefname{originallemma}{Original lemma}{Original lemmas}
\makeatletter
\newlabel{lemma-composition-regular-immersion}{{067Q}{source}{}{sourcepointer-3}{}}
\newlabel{lemma-composition-regular-immersion@cref}{{[originallemma][0][]067Q}{[1][0][]source}{}{}{}}
\makeatother
\crefname{originallemma}{original lemma}{original lemmas}
\Crefname{originallemma}{Original lemma}{Original lemmas}
\makeatletter
\newlabel{lemma-immersion-smooth-into-smooth-regular-immersion}{{067U}{source}{}{sourcepointer-4}{}}
\newlabel{lemma-immersion-smooth-into-smooth-regular-immersion@cref}{{[originallemma][0][]067U}{[1][0][]source}{}{}{}}
\makeatother
\crefname{originallemma}{original lemma}{original lemmas}
\Crefname{originallemma}{Original lemma}{Original lemmas}
\makeatletter
\newlabel{lemma-permanence-regular-immersion}{{068Z}{source}{}{sourcepointer-5}{}}
\newlabel{lemma-permanence-regular-immersion@cref}{{[originallemma][0][]068Z}{[1][0][]source}{}{}{}}
\makeatother
\crefname{originallemma}{original lemma}{original lemmas}
\Crefname{originallemma}{Original lemma}{Original lemmas}
\makeatletter
\newlabel{lemma-regular-quasi-regular-immersion}{{063K}{source}{}{sourcepointer-6}{}}
\newlabel{lemma-regular-quasi-regular-immersion@cref}{{[originallemma][0][]063K}{[1][0][]source}{}{}{}}
\makeatother
\crefname{originallemma}{original lemma}{original lemmas}
\Crefname{originallemma}{Original lemma}{Original lemmas}
\makeatletter
\newlabel{morphisms-lemma-smooth-etale-over-affine-space}{{054L}{source}{}{sourcepointer-7}{}}
\newlabel{morphisms-lemma-smooth-etale-over-affine-space@cref}{{[originallemma][0][]054L}{[1][0][]source}{}{}{}}
\makeatother
\title{Koszul-regularity, H1-regularity and quasi-regularity of an immersion descend through a smooth morphism}
\author{The Stacks Project Authors}
\date{stacks-project:divisors.tex:local-snapshot-20261012; GNU FDL 1.2 or later}
\begin{document}
\maketitle
\noindent\textit{Editorial note.} All delivered mathematical text is the author's own, copied from the pinned Stacks Project chapter file pinned-original-source.tex; no statement or proof step was rewritten, reordered or removed.\par\medskip
\noindent\textit{Editorial note.} Following the multi-claim convention for this source, source\_url is the source-level URL https://stacks.math.columbia.edu/ and the source label lemma-push-regular-immersion-thru-smooth is carried in claim\_key and here rather than in a per-claim tag URL.\par\medskip
\noindent\textit{Editorial note.} The preamble is the chapter's shared Stacks Project preamble, retained as sources/stacks-divisors/retained-preamble.tex: a line-for-line excerpt of the shared Stacks preamble staged with this snapshot. No macro was invented or redefined, and this card adds no preamble text of its own.\par\medskip
\noindent\textit{Editorial note.} Statement span 5740-5753 and proof span 5755-5843 come from the halfA pool card operations/pool-cards/stacks-divisors-halfA.json and were re-read in the pinned chapter file before this card was written.\par\medskip
\noindent\textit{Editorial note.} The statement contains an xymatrix commutative triangle; it is copied verbatim and the retained chapter preamble loads Xy-pic.\par\medskip
\noindent\textit{Editorial note.} The proof contains two author remarks about which reductions fail in which case (for instance the note that adjoining a closed immersion into an affine space does not work in the quasi-regular case) and a '(small detail omitted)'; all are preserved.\par\medskip
\noindent\textit{Editorial note.} 7 cross-references occur inside the counted spans; they are listed in external\_references as navigation pointers only.\par\medskip
\noindent\textit{Allowed standard prerequisites (editorial).}\begin{itemize}
\item The standing equivalence between Koszul-regular, H1-regular and quasi-regular for the theory used here, cited from the same chapter. Cited at proof lines 5755-5756 as used without further mention. Lets the three regularity notions be handled in the same framework.
\item Permanence of regular immersions and regularity of the composite of a regular immersion with a regular immersion, cited from the same chapter. Cited at proof lines 5802-5804 and again at 5826-5828. Closes the first case and lets a closed immersion into an affine space be adjoined without losing regularity.
\item A closed immersion of a smooth scheme into a smooth scheme over the same base is a regular immersion. Cited at proof lines 5818-5819. Produces the regular immersion into an affine space used to reduce the H1-regular case.
\item Factorization of a smooth morphism through an étale morphism to an affine space over the base, together with the formal étale lifting identities for graded pieces of ideals. Cited at proof lines 5834-5836 and 5840-5843. Reduces the quasi-regular case to the affine space case and to the étale case, where the associated graded pieces of the two ideals are identified.
\end{itemize}\medskip
\ifdefined\hypertarget\hypertarget{sourcepointer-1}{}\fi
\noindent\textit{Source reference.} Original lemma 00UR: focused-seeds/AGP/sources/stacks-snapshot-20261011/tags.txt, source lines 1121--1121. Same-chapter reference used in the retained author text at source line(s) 5838. It is supplied only as a navigation pointer so that the retained text compiles; the counted item delivers the author's own statement and the author's complete proof, and no part of the referenced item is copied in or substituted for the counted argument.\par\medskip
\ifdefined\hypertarget\hypertarget{sourcepointer-2}{}\fi
\noindent\textit{Source reference.} Original lemma 0H1D: focused-seeds/AGP/sources/stacks-snapshot-20261011/tags.txt, source lines 21243--21243. Same-chapter reference used in the retained author text at source line(s) 5840. It is supplied only as a navigation pointer so that the retained text compiles; the counted item delivers the author's own statement and the author's complete proof, and no part of the referenced item is copied in or substituted for the counted argument.\par\medskip
\ifdefined\hypertarget\hypertarget{sourcepointer-3}{}\fi
\noindent\textit{Source reference.} Original lemma 067Q: focused-seeds/AGP/sources/stacks-snapshot-20261011/tags.txt, source lines 7979--7979. Same-chapter reference used in the retained author text at source line(s) 5171, 5669, 5817. It is supplied only as a navigation pointer so that the retained text compiles; the counted item delivers the author's own statement and the author's complete proof, and no part of the referenced item is copied in or substituted for the counted argument.\par\medskip
\ifdefined\hypertarget\hypertarget{sourcepointer-4}{}\fi
\noindent\textit{Source reference.} Original lemma 067U: focused-seeds/AGP/sources/stacks-snapshot-20261011/tags.txt, source lines 7983--7983. Same-chapter reference used in the retained author text at source line(s) 5815. It is supplied only as a navigation pointer so that the retained text compiles; the counted item delivers the author's own statement and the author's complete proof, and no part of the referenced item is copied in or substituted for the counted argument.\par\medskip
\ifdefined\hypertarget\hypertarget{sourcepointer-5}{}\fi
\noindent\textit{Source reference.} Original lemma 068Z: focused-seeds/AGP/sources/stacks-snapshot-20261011/tags.txt, source lines 8023--8023. Same-chapter reference used in the retained author text at source line(s) 5807. It is supplied only as a navigation pointer so that the retained text compiles; the counted item delivers the author's own statement and the author's complete proof, and no part of the referenced item is copied in or substituted for the counted argument.\par\medskip
\ifdefined\hypertarget\hypertarget{sourcepointer-6}{}\fi
\noindent\textit{Source reference.} Original lemma 063K: focused-seeds/AGP/sources/stacks-snapshot-20261011/tags.txt, source lines 7834--7834. Same-chapter reference used in the retained author text at source line(s) 4935, 4968, 5297, 5756. It is supplied only as a navigation pointer so that the retained text compiles; the counted item delivers the author's own statement and the author's complete proof, and no part of the referenced item is copied in or substituted for the counted argument.\par\medskip
\ifdefined\hypertarget\hypertarget{sourcepointer-7}{}\fi
\noindent\textit{Source reference.} Original lemma 054L: focused-seeds/AGP/sources/stacks-snapshot-20261011/tags.txt, source lines 6644--6644. Same-chapter reference used in the retained author text at source line(s) 5825. It is supplied only as a navigation pointer so that the retained text compiles; the counted item delivers the author's own statement and the author's complete proof, and no part of the referenced item is copied in or substituted for the counted argument.\par\medskip
\begin{lemma}
\label{lemma-push-regular-immersion-thru-smooth}
Let
$$
\xymatrix{
Y \ar[rd]_j \ar[rr]_i & & X \ar[ld] \\
& S
}
$$
be a commutative diagram of morphisms of schemes.
Assume $X \to S$ smooth and $i$ and $j$ immersions.
If $i$ is a Koszul-regular (resp.\ $H_1$-regular, quasi-regular)
immersion, then so is $j$.
\end{lemma}

\begin{proof}
We will use Lemma \ref{lemma-regular-quasi-regular-immersion}
without further mention.
Let $y \in Y$ be any point. Set $x = i(y)$ and set $s = j(y)$.
It suffices to prove the result after replacing $X$ and $S$ by open
neighbourhoods $U$ and $V$ of $x$ and $s$ and $Y$ by an open neighbourhood
of $y$ in $i^{-1}(U) \cap j^{-1}(V)$.

\medskip\noindent
We first prove the result for $X = \mathbf{A}^n_S$. After replacing $S$
by an affine open $V$ and replacing $Y$ by $j^{-1}(V)$ we may assume that
$j$ is a closed immersions and $S$ is affine. Write $S = \Spec(A)$.
Then $j : Y \to S$ defines an isomorphism of $Y$ to the closed subscheme
$\Spec(A/I)$ for some ideal $I \subset A$. The map
$i : Y = \Spec(A/I) \to
\mathbf{A}^n_S = \Spec(A[x_1, \ldots, x_n])$
corresponds to an $A$-algebra homomorphism
$i^\sharp : A[x_1, \ldots, x_n] \to A/I$.
Choose $a_i \in A$ which map to $i^\sharp(x_i)$ in $A/I$.
Observe that the ideal of the closed immersion $i$ is
$$
J = (x_1 - a_1, \ldots, x_n - a_n) + IA[x_1, \ldots, x_n].
$$
Set $K = (x_1 - a_1, \ldots, x_n - a_n)$. We claim the sequence
$$
0 \to K/KJ \to J/J^2 \to J/(K + J^2) \to 0
$$
is split exact. To see this note that $K/K^2$ is free with basis
$x_i - a_i$ over the ring $A[x_1, \ldots, x_n]/K \cong A$.
Hence $K/KJ$ is free with the same basis over the ring
$A[x_1, \ldots, x_n]/J \cong A/I$. On the other hand, taking derivatives
gives a map
$$
\text{d}_{A[x_1, \ldots, x_n]/A} :
J/J^2
\longrightarrow
\Omega_{A[x_1, \ldots, x_n]/A} \otimes_{A[x_1, \ldots, x_n]}
A[x_1, \ldots, x_n]/J
$$
which maps the generators $x_i - a_i$ to the basis elements $\text{d}x_i$
of the free module on the right. The claim follows. Moreover, note that
$x_1 - a_1, \ldots, x_n - a_n$ is a regular sequence in
$A[x_1, \ldots, x_n]$ with quotient ring
$A[x_1, \ldots, x_n]/(x_1 - a_1, \ldots, x_n - a_n) \cong A$.
Thus we have a factorization
$$
Y \to V(x_1 - a_1, \ldots, x_n - a_n) \to \mathbf{A}^n_S
$$
of our closed immersion $i$ where the composition is
Koszul-regular (resp.\ $H_1$-regular, quasi-regular),
the second arrow is a regular immersion, and the associated conormal
sequence is split. Now the result follows from
Lemma \ref{lemma-permanence-regular-immersion}.

\medskip\noindent
Next, we prove the result holds if $i$ is $H_1$-regular.
Namely, shrinking as in the first paragraph of the proof, we may assume that
$Y$, $X$, and $S$ are affine. In this case we can choose a closed
immersion $h : X \to \mathbf{A}^n_S$ over $S$ for some $n$. Note that
$h$ is a regular immersion by
Lemma \ref{lemma-immersion-smooth-into-smooth-regular-immersion}.
Hence $h \circ i$ is a $H_1$-regular immersion, see
Lemma \ref{lemma-composition-regular-immersion} (note that this step
does not work in the ``quasi-regular case'').
Thus we reduce to the case $X = \mathbf{A}^n_S$ and $S$ affine
we proved above.

\medskip\noindent
Finally, assume $i$ is quasi-regular. After shrinking as in the first
paragraph of the proof, we may use Morphisms, Lemma
\ref{morphisms-lemma-smooth-etale-over-affine-space}
to factor $f$ as $X \to \mathbf{A}^n_S \to S$ where the first
morphism $X \to \mathbf{A}^n_S$ is \'etale. This reduces the problem
to the the two cases (a) $X = \mathbf{A}^n_S$ and (b) $f$ is \'etale.
Case (a) was handled in the second paragraph of the proof.
Case (b) is handled by the next paragraph.

\medskip\noindent
Assume $f$ is \'etale. After shrinking we may assume $X$, $Y$, and $S$
affine $i$ and $j$ closed immersions (small detail omitted).
Say $S = \Spec(A)$, $X = \Spec(B)$ and $Y = \Spec(B/J) = \Spec(A/I)$.
Shrinking further we may assume $J$ is generated by a quasi-regular sequence.
The ring map $A \to B$ is \'etale, hence formally \'etale
(Algebra, Lemma \ref{algebra-lemma-formally-etale-etale}).
Thus $\bigoplus I^n/I^{n + 1} \cong \bigoplus J^n/J^{n + 1}$
by Algebra, Lemma \ref{algebra-lemma-formally-etale-lift-infinitesimal}.
Since $J$ is generated by a quasi-regular sequence, so is $I$.
This finishes the proof.
\end{proof}

\end{document}
